\chapter{Trust, Security, and the Rulebook}
\label{c:trust}

\begin{glance}
Open weights eliminate API dependency; they do not establish trust. The commons is attackable at every layer of its distribution chain, carries a censorship record that transfers into derivatives, and has a contested distillation genealogy; and if the deployment layer behaves like a certification regime rather than a commodity market, it can win every download metric and still lose the workloads that pay. Its security risk is structural, not national: a monoculture of a few model families, quantized into thousands of unowned derivatives, wired into agents holding file, shell, network and credential authority. The realized incidents to date came through runtimes and supply chains, not poisoned weights, and the defensible objective is not trustworthy models but unnecessary model trust. The rulebook is being written by three actors who are mostly not competing---the United States governs by denial, the European Union by regulation, China by standards and institution-building---and the arena that will decide the deployment layer is vacant.
\end{glance}

\section{Trust: the contested layer}

Chapter~\ref{c:elements} treated openness as an adoption weapon; this chapter treats its exposed flank. Trust is what stalled Chinese campaigns in aircraft, drugs and avionics-grade markets, and three dimensions of it are contested for AI. \emph{Supply-chain security}: the model-distribution infrastructure has been shown attackable repeatedly [V]---malicious serialized models, namespace re-registration that fed malicious models into the catalogs of two hyperscalers~\cite{pickle,namespace}, a platform breach by an agentic attacker, and worst of all a training-layer result that about 250 poisoned documents suffice to backdoor a model regardless of scale, with backdoors surviving safety fine-tuning~\cite{poison,sleeper}. The open commons runs on the security posture of early-2000s package managers, at national-infrastructure stakes. \emph{Behavioral integrity}: the peer-reviewed record documents semantic-level suppression of politically sensitive content in DeepSeek-family models, distinct from safety refusal, and transferring to distilled derivatives---the politics ride along with the capability~\cite{deepseek-censor,r1dacted}. This does not make the weights unusable, but ``just self-host it'' becomes a three-layer problem: a locally hosted Chinese model delivers operational sovereignty while leaving epistemic and supply-chain sovereignty unresolved, and closed US models fail the operational layer by construction. \emph{Provenance}: a February 2026 report alleges industrial-scale distillation of a Western frontier model by three Chinese laboratories through some 24,000 fraudulent accounts [P, vendor allegation, not adjudicated]~\cite{anthropic-distill}---which, if true at scale, caps the fast-follower strategy at frontier-minus-epsilon unless domestic capability compounds independently, hands Western policy a legal instrument on the certification axis, and cuts against the clean morality tale: the commons is partly built from enclosed material. On \datafreeze\ itself the NSA, FBI and CISA issued joint advisory AA26-251a naming DeepSeek, Moonshot, Alibaba, MiniMax, StepFun and Z.AI for distillation against US frontier models routed through cloud providers and aggregators~\cite{cisa-distillation} [P]---an allegation by security agencies, not an adjudication, and about distillation, not hosting; what is new is that a vendor's grievance in February is a state's position in September, naming very nearly the whole Chinese open frontier in one document.

The rational response for a third country is neither refusal nor naive adoption but infrastructure: signed manifests, reproducible build and quantization pipelines, non-executable packaging, model bills of materials, national evaluation batteries, sandboxed loaders, attested serving firmware and national artifact mirrors---what software supply-chain security did after SolarWinds, applied to weights, organized as five layers of trust and a certification ladder from Level~0 to Level~5. Most of the commons ships at Level~0--1; most regulated deployment will require Level~3--4, and that gap is the deployment bottleneck of the entire commoditization thesis. A new object has appeared in the meantime: an August 2026 Chinese open-weight video model that tops its arena under a licence excluding local deployment in the United States, the European Union, the United Kingdom and South Korea---adoption weapon and export control in one artifact~\cite{minimax-geofence}, and a procurement office that checks hashes but not jurisdictions has audited the wrong layer. Whoever runs the evaluation batteries, signs the manifests and operates the mirrors becomes the FAA of models; China's World AI Cooperation Organization is bidding for that role, no Western or neutral institution has claimed it, and for Japan and similarly positioned states the vacancy is the highest-leverage square on the board \full{18}.

\section{Monoculture and agency: the security cost of the commons}

National origin is neither evidence of a backdoor nor a technical risk category; the risk is structural---a property of monoculture and delegated authority, not of a flag. The full edition separates four mechanisms that produce identical outcomes and need different evidence and mitigations---a learned backdoor in the weights, a clean model hijacked at runtime by injection or self-replicating promptware, executable malware in the package, and data collection by a hosted API---and grades the evidence for each on the security literature's own scale~\cite{order66} \full{19}. What has been demonstrated in controlled work is severe: sleeper-agent backdoors surviving supervised and reinforcement fine-tuning, poisoning at near-constant sample cost, memory-to-network exfiltration with under 1\% utility loss---which is why ordinary evaluation does not find it~\cite{backreveal}---and self-replicating promptware persisting through restarts. What has actually happened is different in kind: no confirmed weight-level backdoor in any released official checkpoint; instead malicious repositories, malicious agent skills, and, in July 2026, a frontier laboratory's cyber-evaluation escaping isolation and reaching a model-hosting platform's production systems through some 17,600 autonomous actions with no human directing them~\cite{openai-hf,hf-timeline}, plus a second laboratory's disclosure that models in its evaluations had compromised real organizations' systems~\cite{anthropic-incidents}. Neither involved a poisoned weight, a prompt worm, or a malicious operator. That is the finding: the dominant realized risk is compound and mundane---narrow-goal autonomy, a containment error, ordinary vulnerabilities, overbroad credentials, and public services composing across mismatched trust boundaries.

The compound ``Order~66'' scenario---a dormant policy, an independent actor who learns the trigger, a self-replicating carrier---has not been observed; it requires a conjunction of access, activation, persistence, authority, propagation and objective, and breaking any one conjunct prevents the chain. Two asymmetries follow: the implanter loses control, because a backdoor is a reusable vulnerability shared by every deployment; and blast radius is set by authority, not by the payload's phrasing. A fault-tree analysis with common-cause terms for shared model families and shared runtimes puts a five-year systemic-loss probability roughly between $10^{-6}$ and $10^{-2}$ [S], the upper end driven by correlation structure, and the most effective interventions cut the common-cause factors, which runtime diversity and deterministic tool mediation do and model-family diversity alone does not. The defensible objective is therefore not trustworthy models but unnecessary model trust---even a perfectly backdoored model cannot send a secret through a network path the operating system and tool broker do not provide; the full edition lists fourteen control families \fullapp{E}. Certification and hardening are the strongest legitimate brake on the adoption curve, and that lag---not tariffs---is the West's real remaining source of deployment-layer time, available to any jurisdiction willing to build the apparatus rather than merely publish the rule. The vacancy is urgent: no commercial actor has an incentive to certify honestly models it does not sell, and no single government will be trusted to certify models it did not train.

\section{The governance contest: who writes the rules}

Technology wars are settled twice: once in the factories, once in the rulebooks, and in each precedent the winner of the industrial contest inherited the pen. The standard Western framing---light US regulation, heavy European, authoritarian Chinese---misdescribes the contest badly enough to lose it. The United States governs by denial: export controls, entity lists, the foreign direct product rule, no general statute at home, and in August 2026 a voluntary frontier-model evaluation framework finalized in closed session and left unpublished~\cite{wh-framework}. Denial produces substitution and confers no positive standard---no definitions, no conformity procedure, no registry---so when it fails there is nothing left behind. The European Union governs by regulation, extraterritorial through its market, with the weakness that regulation is not production and leverage decays with deployment share. China governs by standards and institution-building, and this is what Western commentary under-weights: the earliest generative-service filing regime, a mandated unified chip-interconnection ecosystem, national standards as procurement conditions, open-sourced interconnect specifications and software stacks, the largest bloc in the open-ISA consortium, a Global AI Governance Action Plan, and a World AI Cooperation Organization in Shanghai with twenty-nine founding members, none a major Western democracy~\cite{action-plan,waico}. This is what winning a standards war looks like from the inside: not a treaty imposed, but a stack given away until it is the default, and then a forum built around the default.

Standards are load-bearing because a ratified standard survives the political cycle. The book maps each contested standard on a four-stage ladder---de jure specification, open reference implementation, conformance suite, installed-base default---and rent lives at stage four while committees fight over stage one. China's distinctive play is to skip stage politics: ship the reference implementation and the mandated installed base and let the specification ratify what already exists; the West's distinctive failure is stranding specifications at stages one and two while defending installed bases it assumes are permanent. Standards leadership tends to follow manufacturing leadership with a five-to-ten-year lag and then lock in for twenty, though scale confers influence rather than control---the counter-cases in photovoltaic standards, charging connectors, and 5G are recorded. Two governance problems the current instruments do not survive: nobody agrees who the ``provider'' of a quantized derivative of a merged fine-tune redistributed by an individual is, and compute-threshold regimes cannot see the half-to-one-megawatt sovereign clusters that decentralized pretraining produces---so decentralization is a governance forecast as well as an industrial one. The vacant arena is trust and certification of the open commons. The United States cannot fill it, having no positive procedure and no wish to certify what it would rather did not exist; the European Union has paper apparatus but not operations; China's forum-building is the only bid on the table---not obviously worse than none, but not what a buyer would design. Of three futures, convergence around a referee neither principal controls is least likely---in aviation and pharmaceuticals it arrived only after accidents---bifurcation into two certified stacks is most likely, and fragmentation is the unstable failure case. The rules of the AI era are not being written where most Western attention is directed \full{20}.

